\subsection{Data Preprocessing}
\label{sec:data-preprocessing}
% Rewritten 2026-11-01 in the order of the actual pipeline (see notes/2_BOARD.md, Methods M2); polished 2026-11-02 (P5).

\noindent \textbf{Candidate peak--gene pairs.} CIMA mapped eQTLs for each gene using all SNPs within 1 Mb of its transcription start site (TSS), and caQTLs for each chromatin accessibility peak using all SNPs within 1 Mb of the peak midpoint \cite{yin2026cima}. We considered all genes that had CIMA eQTL results and were present in the pseudobulk RNA data, together with all peaks in the CIMA ATAC data. The TSS of each gene was recovered from the eQTL results as the variant position minus the reported distance to the gene start. A peak--gene pair was retained as a candidate if the distance between the TSS and the peak center was at most 1 Mb.

\noindent \textbf{Instrument selection.} To satisfy the relevance assumption, we retained only caQTL variants with nominal $P<10^{-8}$ in each cell type. We then performed LD clumping separately for each peak--gene pair with PLINK (v1.90b7.2) through \texttt{ieugwasr::ld\_clump\_local} (ieugwasr 1.1.0), using an $r^2$ threshold of 0.99, a window of 20,000 kb that covers the entire \textit{cis} region, index variants ranked by caQTL $P$ value, and \texttt{clump\_p} = 1. The purpose of clumping was not to obtain independent SNPs but to avoid near-perfect collinearity among highly correlated instruments; the remaining correlation was modeled by GLS (\nameref{sec:gls-ivw}). We kept pairs for which the caQTL and eQTL data shared at least three variants and for which both the gene and the peak were present in the pseudobulk and WNN data (\nameref{sec:pseudobulk-correlations} and \nameref{sec:wnn-correlations}). At least three instruments were required to remain after each subsequent filtering step.

\noindent \textbf{Allele harmonization and duplicate instruments.} In both the caQTL and eQTL summary statistics, the alternative (ALT) allele was the effect allele. CIMA had harmonized the two sets of summary statistics, and we verified that their effect and non-effect alleles agreed; variants with discordant alleles would have been removed. Some SNPs had nearly identical caQTL and eQTL associations. A SNP was removed if both its caQTL and eQTL slopes differed by at most $10^{-3}$ from those of a previously retained SNP. Because SNPs were processed in variant-ID order, the first SNP of each such set was kept.

\noindent \textbf{LD matrix.} The signed LD correlation matrix $\boldsymbol R$ among the retained instruments was computed with PLINK (v1.90b7.2) using \texttt{--r square --keep-allele-order --write-snplist}, so that the correlations refer to the same effect alleles as the summary statistics. Variants absent from the reference panel, as well as rows and columns with undefined (NaN) correlations, were removed. The covariance matrix used in GLS was $\boldsymbol{\Sigma}=\operatorname{diag}(s_{Y})\,\boldsymbol R\,\operatorname{diag}(s_{Y})$ (\nameref{sec:gls-ivw}). Instruments were then oriented so that all caQTL effects were positive, and $\boldsymbol R$ was reoriented accordingly (\nameref{sec:orientation}). For both clumping and LD estimation, the reference panel was the East Asian (EAS) subset ($n=585$) of the 1000 Genomes Project GRCh38 release (3,202 individuals) [CITE: 1000 Genomes GRCh38 high-coverage release]; chromosome X was coded as 23.
\TODO{Confirm that the two EAS panel files used for clumping and for the LD matrix contain the same variant set (A1-09).}
